%! Characteristics of Linear Functions — Unit 1, Lesson 1.0 — Guided Notes (Answer Key)
% Mirrors notes/main.tex exactly; every \blank{} becomes an \ans{} and every work block is
% authored byte-identically, so blank and key paginate alike.
\documentclass[10pt]{article}
\usepackage{algebra2-article}
\usepackage{algebra2-key}   % pulls in -boxes; do NOT also load -boxes
\newcommand{\TallMath}[1]{$\displaystyle #1\rule[-1.4em]{0pt}{3.2em}$}
% Temperature graph for the Individual Practice: hours h on 0..7, temperature T on -8..8.
\newcommand{\tempgrid}{%
  \draw[step=1, linegray!40, very thin] (0,-8) grid (7,8);
  \draw[->, semithick, charcoal] (0,0)--(7.6,0) node[right, font=\scriptsize]{$h$};
  \draw[->, semithick, charcoal] (0,-8.6)--(0,8.8) node[above, font=\scriptsize]{$T$ ($^\circ$C)};
  \foreach \x in {1,...,7} \node[charcoal, font=\tiny, below] at (\x,0){\x};
  \foreach \y in {-8,-6,-4,-2,2,4,6,8} \node[charcoal, font=\tiny, left] at (0,\y){\y};}
% Coordinate grid: {xmin}{xmax}{ymin}{ymax}
\newcommand{\lgrid}[4]{%
  \draw[step=1, linegray!40, very thin] (#1,#3) grid (#2,#4);
  \draw[->, semithick, charcoal] (#1,0)--(#2,0) node[below right, font=\scriptsize]{$x$};
  \draw[->, semithick, charcoal] (0,#3)--(0,#4) node[left, font=\scriptsize]{$y$};}
% Vocab answer for the key: bold forest term, then the filled definition on a write-line.
% The leading and trailing \par are required: \ansline ends with \dotfill but does NOT end the
% paragraph, so without them each term label gets pulled onto the previous answer's line.
\newcommand{\vocabans}[2]{%
  \par\noindent\textbf{\textcolor{forest}{#1:}}\\[1pt]\ansline{#2}\par}

\begin{document}

\pageheader{Unit 1, Lesson 1.0}{Guided Notes --- Answer Key}

\vspace{0.10in}

\begin{objectivebox}
\textbf{By the end of this lesson, I will be able to\dots}
\begin{itemize}\setlength{\itemsep}{2pt}
  \item recognise a \emph{linear} function by its \emph{constant rate of change}, and find that rate
        (the \emph{slope}) from a graph, a table, or two points;
  \item find the \emph{$y$-intercept} and the \emph{zero} ($x$-intercept) from a graph, a table, or a rule;
  \item say where a line is \emph{increasing}, \emph{decreasing}, or \emph{constant}, and where it is
        \emph{positive} or \emph{negative};
  \item state the \emph{domain} and \emph{range} of a line --- as pure math, and as a story allows.
\end{itemize}
\end{objectivebox}

\vspace{0.05in}

\begin{vocabbox}
\small
Fill in each term as we build it together.
\par\vspace{2pt}   % \par is required: \vocabans's \noindent is a no-op mid-paragraph
\vocabans{Linear function}{a function with a constant rate of change; its graph is a straight line, and it can be written $f(x)=mx+b$}
\vocabans{Slope (rate of change)}{$m=\Delta y/\Delta x=(y_2-y_1)/(x_2-x_1)$ --- the change in output per unit change in input}
\vocabans{$y$-intercept}{the point $(0,b)$ where the graph crosses the $y$-axis; the output when $x=0$ --- the ``starting value''}
\vocabans{$x$-intercept (zero)}{the point $(a,0)$ where the graph crosses the $x$-axis; the input $a$ with $f(a)=0$}
\vocabans{Increasing / decreasing / constant}{rising ($m>0$) / falling ($m<0$) / flat ($m=0$), read left to right}
\vocabans{Positive / negative interval}{the $x$-values where the graph sits above ($f(x)>0$) or below ($f(x)<0$) the $x$-axis}
\end{vocabbox}

\begin{hookbox}
At 6:00 a.m.\ Saturday it was $-6\,^\circ$C in Blacksburg, and the temperature rose a steady
$2\,^\circ$C every hour. Let $h$ be the hours after 6:00 a.m.
\begin{itemize}\setlength{\itemsep}{1pt}
  \item $T(0)=\ans{$-6$}$ \quad $T(1)=\ans{$-4$}$ \quad $T(2)=\ans{$-2$}$ \quad $T(4)=\ans{$2$}$
  \item The temperature reaches $0\,^\circ$ at $h=\ans{$3$}$ --- and \emph{that} is the hour the
        roads stop being icy.
\end{itemize}
Two numbers tell the whole story: \emph{where it starts} and \emph{how fast it changes}. Today we
name every feature of a line, and read each one three ways --- off a graph, a table, and a rule.
\end{hookbox}

\boxguard[24]
\begin{notesbox}{1. Linear Means the Rate of Change Is Constant}
Fill in the Saturday table, then the row of changes underneath it.
\begin{center}
\renewcommand{\arraystretch}{1.3}
\begin{tabular}{|c|c|c|c|c|c|}
\hline
$h$ & $0$ & $1$ & $2$ & $3$ & $4$ \\ \hline
$T$ & $-6$ & $-4$ & \ans{$-2$} & \ans{$0$} & \ans{$2$} \\ \hline
change in $T$ & --- & $+2$ & \ans{$+2$} & \ans{$+2$} & \ans{$+2$} \\ \hline
\end{tabular}
\end{center}
The change is the \ans{same $+2$} every step --- that is what makes the function
\textbf{linear} and the graph \ans{straight}. We call that constant the \textbf{slope}.
\begin{center}
\begin{minipage}[c]{0.46\linewidth}
\[ m \;=\; \frac{\Delta y}{\Delta x} \;=\; \frac{\text{rise}}{\text{run}}
        \;=\; \frac{y_2-y_1}{x_2-x_1} \]
Every linear function can be written
\[ f(x) \;=\; mx + b. \]
\end{minipage}\hfill
\begin{minipage}[c]{0.48\linewidth}\centering
\begin{tikzpicture}[scale=0.44]
  \lgrid{-2}{6}{-6}{6}
  \draw[very thick, forest] (-1,-6)--(5,6);
  \draw[semithick, navy, dashed] (2,0)--(3,0)--(3,2);
  \node[navy, font=\tiny] at (2.5,-0.8){run $1$};
  \node[navy, font=\tiny, right] at (3.1,1){rise $2$};
  \fill[charcoal] (0,-4) circle (3.4pt);
  \fill[charcoal] (2,0) circle (3.4pt);
\end{tikzpicture}\\[1pt]
{\footnotesize $f(x)=2x-4$}
\end{minipage}
\end{center}
\textbf{Slope from any two points.} The line above passes through $(1,-2)$ and $(3,2)$.
\begin{work}
  m &= \frac{2 - (-2)}{3 - 1} \\
    &= \frac{4}{2} \\
    &= 2
\end{work}
\emph{Any} two points give the same answer --- because the rate of change is \ans{constant}.
\end{notesbox}

\boxguard[22]
\begin{notesbox}{2. Where the Graph Meets the Axes}
\begin{center}
\renewcommand{\arraystretch}{1.4}
\begin{tabularx}{\linewidth}{@{}l X X@{}}
\textbf{Feature} & \textbf{How to find it} & \textbf{What it means in a story} \\
\hline
$y$-intercept $(0,b)$ & set $x=\ans{$0$}$ & the \ans{starting} value \\
$x$-intercept, or \textbf{zero}, $(a,0)$ & solve $f(x)=\ans{$0$}$ & the moment the quantity \ans{runs out} \\
\end{tabularx}
\end{center}
\textbf{Worked: $f(x)=2x-4$.} \; Its $y$-intercept is $(0,\ans{$-4$})$ --- read it straight
off the rule. For the zero, solve:
\begin{work}
  2x - 4 &= 0 \\
      2x &= 4 \\
       x &= 2
\end{work}
so the zero is at $(\ans{$2$},\,\ans{$0$})$. Check it against the graph above.
\end{notesbox}

\boxguard[26]
\begin{notesbox}{3. Which Way It Moves --- and Where It Sits}
These are \emph{two different questions}. Keep them apart.
\begin{center}
\renewcommand{\arraystretch}{1.35}
\begin{tabularx}{\linewidth}{@{}X|X@{}}
\textbf{Which way it moves} (read left to right) & \textbf{Where it sits} (above or below the $x$-axis) \\
\hline
$m>0$: the line is \ans{increasing} & $f(x)>0$ where the graph is \ans{above} the axis \\
$m<0$: the line is \ans{decreasing} & $f(x)<0$ where the graph is \ans{below} the axis \\
$m=0$: the line is \ans{constant} & the sign can only switch at a \ans{zero} \\
\end{tabularx}
\end{center}
\textbf{For $f(x)=2x-4$:} it is \ans{increasing}; it is positive when $x \ans{$>2$}$
and negative when $x \ans{$<2$}$.

\vspace{3pt}
\begin{tcolorbox}[colback=goldbg, colframe=goldacc, boxrule=0.8pt, arc=1.5mm,
    left=2.5mm, right=2.5mm, top=1.5mm, bottom=1.5mm]
\textbf{\textcolor{goldacc}{Careful: ``positive'' is not ``increasing.''}} \;
Sunday's temperature was $g(h)=6-2h$. At $h=1$ it is $g(1)=\ans{$4$}$ degrees --- a
\ans{positive} number --- and the temperature is \ans{falling}. Both at once.
\emph{Positive} says where the graph \textbf{sits}; \emph{increasing} says which way it
\textbf{moves}.
\end{tcolorbox}
\end{notesbox}

\boxguard[20]
\begin{notesbox}{4. Domain and Range --- the Math, and the Story}
\begin{itemize}[leftmargin=1.5em, itemsep=3pt]
  \item A line that is \emph{not} horizontal: domain \ans{all reals}, range \ans{all reals}.
        The rule accepts every input and produces every output.
  \item A \textbf{horizontal} line $c(x)=b$ is the exception: slope \ans{$0$}, neither
        increasing nor decreasing, and range \ans{$\{b\}$} --- a single value. It has no zero
        at all unless $b=\ans{$0$}$, and then \emph{every} input is one.
  \item In a \textbf{story}, the domain shrinks to the inputs that \ans{make sense}. Saturday's
        rule works for $h=-100$; Saturday morning does not.
\end{itemize}
\end{notesbox}

\begin{practicebox}
Everything at once, on a new line: $h(x) = -2x + 6$.
\begin{enumerate}[leftmargin=1.7em, itemsep=3pt]
  \item Slope \ans{$-2$}; \quad $y$-intercept \ans{$(0,6)$}; \quad increasing or decreasing?
        \ans{decreasing}
  \item Find the zero.
        \begin{work}
          -2x + 6 &= 0 \\
             -2x  &= -6 \\
               x  &= 3
        \end{work}
        Zero: $(\ans{$3$},\,\ans{$0$})$.
  \item $h(x)>0$ when $x \ans{$<3$}$; \quad $h(x)<0$ when $x \ans{$>3$}$.
  \item Domain \ans{all reals}; \quad range \ans{all reals}.
\end{enumerate}
\end{practicebox}

\boxguard[22]
\begin{scenariobox}[Individual Practice --- On Your Own]{navy}
Three problems, quietly and on your own. Keep the two questions apart --- \emph{which way it
moves} and \emph{where it sits} --- and for every answer, be able to point to it on the graph.
\begin{enumerate}[leftmargin=1.9em, itemsep=8pt]
  \item \textbf{Sunday evening.} At 4:00 p.m.\ on Sunday it was $6\,^\circ$C, and the temperature
        \emph{fell} a steady $2\,^\circ$C every hour; $h$ counts the hours after 4:00 p.m. \\[2pt]
    \begin{minipage}[c]{0.60\linewidth}
    Rule: $T(h)=$ \ans{$6-2h$} \quad slope \ans{$-2$}; \; \ans{decreasing} (increasing /
    decreasing) \\[3pt]
    $y$-intercept \ans{$(0,6)$}, meaning \ans{$6^\circ$ at 4:00 p.m.} \\[3pt]
    Zero: solve $T(h)=0$, then circle it on the graph.
    \end{minipage}\hfill
    \begin{minipage}[c]{0.36\linewidth}\centering
    \begin{tikzpicture}[x=0.40cm, y=0.18cm]
      \tempgrid
      \draw[very thick, navy] (0,6)--(7,-8);
      \fill[charcoal] (0,6) circle (2.8pt);
    \end{tikzpicture}
    \end{minipage}
    \begin{work}
      6 - 2h &= 0 \\
         -2h &= -6 \\
           h &= 3
    \end{work}
    $T$ is positive when $h$ \ans{$<3$} and negative when $h$ \ans{$>3$}. \\[3pt]
    \textbf{The crux.} At 5:00 p.m.\ ($h=1$) the temperature is \ans{$4$} $^\circ$C --- a
    \ans{positive} number --- and it is \ans{falling} (rising / falling). A classmate says ``it
    can't be positive, it's decreasing.'' What are they mixing up?
    \par\ansline{\emph{Positive} says where the graph sits --- above the axis; \emph{decreasing} says which way it moves.}
    \ansline{Both are true at $h=1$: $T=4>0$, and the line is heading down.}

  \item \textbf{Side by side.} Saturday's line $T=2h-6$ and Sunday's $T=6-2h$ share the zero
        $h=3$. Saturday is negative \ans{before} $h=3$ and positive \ans{after} it; Sunday is
        the other way round. Their slopes are \ans{$2$} and \ans{$-2$}. Why can neither
        line change sign anywhere \emph{except} at $h=3$?
    \par\ansline{A line's sign can only switch where it crosses the $h$-axis --- at a zero.}
    \ansline{Each line has exactly one zero, at $h=3$, so it changes sign exactly once, there.}

  \item \textbf{Model it.} A car-wash gift card starts with \$40 and each wash costs \$8, so the
        balance after $w$ washes is $B(w)=40-8w$. \\[2pt]
    (a) The slope $-8$ means \ans{each wash takes \$8 off the balance}; the $y$-intercept $(0,40)$ means \ans{the card starts with \$40} \\[3pt]
    (b) Find the zero.
    \begin{work}
      40 - 8w &= 0 \\
          -8w &= -40 \\
            w &= 5
    \end{work}
    In the story, the zero means: \ans{after $5$ washes the card is empty} \\[3pt]
    (c) Which inputs $w$ make sense? \ans{$w=0,1,2,3,4,5$} \quad (d) If each wash cost \$10 instead, would
        the card run out \emph{sooner} or \emph{later}? \ans{sooner} \; What changes on the graph
        --- where the line \emph{starts}, how \emph{steep} it is, or both? \ans{only the steepness --- it still starts at \$40}
\end{enumerate}
\end{scenariobox}

\end{document}
